\documentclass{ximera}

\title{Werkbundel: Standaardvorm}
\author{Daan Lipkens}

\begin{document}
\begin{abstract}
    Inleiding tot het SI-stelsel, eenheden omzetten, beduidende cijfers en de wetenschappelijke notatie.
\end{abstract}
\maketitle

\section*{Inleiding}\label{hoofdstuk:standaardvorm}
Op 21 februari 2023 schreef \textit{The Jerusalem Post} over een meteoor die boven Texas ontplofte. Volgens de krant was hij ``zo groot als een corgi'' (figuur~\ref{fig:corgi}) en zo zwaar als vier babyolifanten.

% AFBEELDING: corgi
% bron: https://commons.wikimedia.org/wiki/File:Welsh_Pembroke_Corgi.jpg
\begin{figure}[ht!]
    \captionsetup{width=0.85\textwidth}
    \begin{center}
        \includegraphics[width=0.45\textwidth]{pic5/appendix/corgi.jpg}
        \fotobron{Dog breed facts}{CC BY-SA 4.0}{https://commons.wikimedia.org/wiki/File:Welsh_Pembroke_Corgi.jpg}
    \end{center}
    \caption{Een Pembroke Welsh corgi. Maar hoe groot is ``zo groot als een corgi'' nu precies?}
    \label{fig:corgi}
\end{figure}

Maar wat is nu de omvang van die meteoor? Een pasgeboren olifant heeft bijvoorbeeld typisch een massa van ongeveer $90 \, \mathrm{kg}$ -- $136 \, \mathrm{kg}$, terwijl een olifant die nog maar 6 maanden oud is al een massa rond de $250 \, \mathrm{kg}$ heeft. Dit is een groot verschil! En wat met de grootte van een corgi? Gaat het over een mannetje of een vrouwtje? Een puppy of een volwassen hond? Je voelt op deze manier aan dat er heel wat nood is aan een vast metrisch systeem zodat lengtematen, massa's ... eenduidig vastgelegd zijn. Daarom heeft men het \textit{Système international d'unités} (of kortweg SI-stelsel) in het leven geroepen.

\end{document}